\section{Results: the AME benchmark (diversity at matched fidelity)}
\label{sec:res:ame}

\plan{\textbf{Claim.} On an established public benchmark, \cleo{} produces
libraries with the sequence coverage of high-temperature LigandMPNN and the
pass rate of low-temperature LigandMPNN --- a combination no single temperature
achieves. \textbf{Why this section exists.} The three application sections rest
on experimental assays that take months and cannot be re-run by a referee. This
one is fully computational, uses a published benchmark and a published
deposit, and can be reproduced end to end. It is the section that makes the
method credible before any wet-lab claim is evaluated. Status:
\statusscaffold{} --- the \cleo{} arm is measured, the temperature sweep is
not yet run.}

\subsection*{Block 1 --- Setup and the unit of comparison}

We evaluate on the atomized motif enforcement (AME) benchmark introduced with
RFdiffusion2, comprising 41 catalytic sites scaffolded into 100 backbones
each. The published protocol assigns 8 LigandMPNN sequences per backbone and
predicts each with 5 Chai models; a prediction succeeds when the motif
all-atom RMSD to the design is below \SI{1.5}{\angstrom} and no backbone atom
clashes with the ligand.

\plan{The aggregation rule matters more than it looks and is easy to get wrong.
We verified it against the deposit rather than assuming: the per-site number
reported in \texttt{AME\_benchmark\_summary.csv} is the fraction of the 100
backbones with \emph{at least one} passing prediction out of the 40, matching
to zero error across all 41 sites and both methods. It is a best-of-40, not a
mean over predictions. Scoring one prediction per sequence and comparing rates
therefore understates a comparable method by roughly threefold --- in the
deposit, only 36\,\% of the sequences that pass on some Chai model pass on all
five. Every number below is on their unit.}

\begin{planlist}
  \item Pilot backbones: three drawn from the \texttt{pass} class, at
        32/40, 28/40 and 10/40 passing predictions in the deposit.
  \item Our predictor is AlphaFold3 rather than Chai; best-of-$N$ uses 5
        diffusion samples per sequence.
  \item \textbf{Caveat to resolve before submission.} RFdiffusion2 scores
        against the unidealized, packed design; we score against the idealized
        backbone. This is the leading explanation for two of the three pilot
        backbones failing in our pipeline at the published temperature despite
        passing in the deposit. Until it is closed, only the backbone that
        reproduces (M0097) supports a head-to-head claim.
\end{planlist}
\todo{ame --- close the idealized/unidealized reference-structure discrepancy}

\subsection*{Block 2 --- Low-temperature sampling returns one solution,
sampled repeatedly}

The incumbent protocol samples LigandMPNN at low temperature. On M0097 this
yields a high pass rate --- 5 of 8 sequences --- but the five passing designs
are \SI{79}{\percent} identical to one another and \SI{87}{\percent} identical
to their own consensus. Under complete-linkage clustering they collapse to a
\emph{single} cluster at every threshold down to \SI{50}{\percent} identity.
The library is one solution measured five times, and the folding budget spent
distinguishing its members buys nothing.

Raising the temperature removes the redundancy and the signal together. At
$T=1.0$ the same backbone yields sequences at \SI{31}{\percent} mutual identity
and \emph{zero} passing designs out of eight. Diversity alone is not the
objective; a library that cannot be filtered is not a library.

\subsection*{Block 3 --- \cleo{} occupies the corner neither temperature
reaches}

\plan{Numbers below are measured. Best-of-5 AF3 diffusion samples, 8-sequence
draws to match the benchmark unit, M0097.}

\begin{center}
\begin{tabular}{@{}lccc@{}}
\toprule
\textbf{Arm} & \textbf{Passing / 8} & \textbf{Distinct clusters / 8} &
\textbf{Coverage vs.\ $T=1.0$} \\
\midrule
LigandMPNN $T=0.1$ & 5.0 & 1.0 & \SI{48}{\percent} \\
LigandMPNN $T=1.0$ & 0.0 & --- & \SI{100}{\percent} \\
\cleo{}-GRPO       & 3.1 & 3.1 & \SI{96}{\percent} \\
\bottomrule
\end{tabular}
\end{center}

\plan{Coverage is the count of variable sequence positions at matched $n=8$,
as a fraction of the $T=1.0$ value; the same ordering holds for per-position
amino-acid count and entropy, where \cleo{} reaches \SIrange{72}{81}{\percent}
and \SIrange{68}{78}{\percent} of $T=1.0$ across the three backbones. Report
positional coverage rather than entropy in the main text and both in the
supplement --- ``comparable coverage'' is defensible at
\SIrange{90}{96}{\percent}, ``same coverage'' is not.}

Two features of this table carry the claim. First, \cleo{} explores
\SI{96}{\percent} of the sequence positions that $T=1.0$ explores while
converting an unfilterable library into a filterable one. Second, for
\cleo{} the passing count and the cluster count coincide: every passing design
is its own solution, whereas the low-temperature arm returns five designs and
one solution. Per unit of folding compute, \cleo{} returns roughly threefold
more independent solutions.

The two arms are not exploring the same region. Every passing \cleo{} design
lies at least 104 mutations --- \SI{58}{\percent} of a 180-residue protein ---
from the nearest passing low-temperature design, and the best \cleo{} design
scores marginally better than the best low-temperature design
(\SI{1.14}{\angstrom} vs.\ \SI{1.16}{\angstrom}). The neighborhood that
LigandMPNN's likelihood concentrates on is one satisfying region among many,
and there is no evidence it is the best one.

\subsection*{Block 4 --- The training trajectory is the library}

\plan{This is the reframing that removes a hyperparameter. Rather than adding
an explicit diversity term to the reward, we treat every rollout sampled during
training as a library candidate. Those sequences were already folded and scored
to compute the reward, so the candidate pool is free: the library is a
byproduct of training, not an additional expense.}

Across 200 GRPO steps at 16 sequences per step, M0097 accumulates 3200 scored
designs, of which 97 pass. Those 97 are \SI{33}{\percent} identical on average
and resolve into 97 distinct clusters at \SI{90}{\percent} identity --- no two
passing designs are near-duplicates. The count is still rising at the end of
training rather than saturating.

\begin{center}
\begin{tabular}{@{}lccc@{}}
\toprule
\textbf{Cumulative folds} & \textbf{Unique passing} &
\textbf{Clusters @\,\SI{90}{\percent} id} & \textbf{Distinct residue identities} \\
\midrule
400  & 14 & 14 & 848 \\
800  & 51 & 51 & 1395 \\
1600 & 85 & 85 & 1887 \\
2400 & 90 & 90 & 1962 \\
3200 & 97 & 97 & 2121 \\
\bottomrule
\end{tabular}
\end{center}

\plan{Honest reading of the deceleration: the curve flattens between 1600 and
2400 folds, but the cause is policy decay rather than exhaustion of sequence
space. Batch diversity bottoms out at the performance peak (steps 36--45) and
then rises while pass rate falls --- with \texttt{kl\_weight} at zero and a
fixed-bounds reward whose within-batch spread collapses once median RMSD
reaches \SIrange{2}{3}{\angstrom}, the policy drifts off its optimum rather
than continuing to search. Rank-normalised rewards and early stopping are
expected to keep the curve climbing; this is the ablation to run before the
figure is final.}

\plan{Pool numbers in this block are single-sample AF3 scores, because that is
what training computed. They are \emph{not} interchangeable with the
best-of-5 numbers in Block 3 and must not be placed on the same axis.}

\subsection*{Block 4b --- Where in sequence space the passing designs sit}

\plan{Generated by \texttt{paper/figures/ame\_pca.py}, which writes two
figures because one projection cannot honestly carry both claims.

\texttt{ame\_pca.svg} projects all arms into a single shared PCA of one-hot
encoded sequences, with \cleo{} restricted to its peak window so that three
\emph{policies} are compared at one point in training. The low-temperature arm
appears as a tight knot displaced far along PC1; \cleo{} occupies a broad cloud
whose passing designs are spread throughout it rather than concentrated at one
edge.

\texttt{ame\_pca\_drift.svg} is the control that keeps the first figure honest.
PC1 of the \emph{pooled} trajectory correlates with training step at
$r = +0.92$ to $+0.94$ on all three backbones, so most of the spread along that
axis is drift over training, not diversity available from any single policy.
That distinction is load-bearing: drift-derived spread is genuine library
diversity when the library is assembled from the whole run, which is how we
assemble ours (Block 4), but it is \emph{not} evidence that one policy is
diverse. Reporting the pooled projection alone would conflate the two and
overstate the result. Fit the shared projection across arms --- a per-arm PCA
places every cloud at the origin and destroys the comparison.

Two limitations to state in the caption. Explained variance is low
(PC1 \SIrange{5}{9}{\percent}), which is expected for one-hot protein sequence
and means the projection is a visual aid, not a metric --- all quantitative
diversity claims use Hamming distance and cluster counts directly. And the
baseline arms currently hold 8 sequences each, so their apparent compactness is
partly small-$n$; the temperature sweep in Block 5 is what makes the clouds
comparable.}
\todo{ame --- rebuild PCA panels once the temperature sweep supplies matched $n$}

\subsection*{Block 5 --- The temperature frontier}

\plan{\textbf{The experiment that decides the section.} Sample the base
LigandMPNN at $T \in \{0.1, 0.2, 0.3, 0.5, 0.7, 1.0\}$, fold a budget matched
to the \cleo{} trajectory pool, and apply the identical library-selection
procedure to every arm --- greedy maximum-diversity selection over the passing
subset. Plot pass rate against passing-set diversity. The claim is that the
baseline traces a frontier along which fidelity and coverage trade off
one-for-one, and that \cleo{} sits above it.

\textbf{This section fails if} the sweep reaches \cleo{}'s coverage at
\cleo{}'s pass rate at some intermediate temperature. Then there is no frontier
shift, only a reparameterisation, and the honest conclusion is that tuning
temperature was sufficient all along. $T=0.2$ and $T=0.3$ are the dangerous
points and must be run --- reporting only $T=0.1$ and $T=1.0$, the two
endpoints we happen to have, would be selecting the comparison that flatters
us.

\textbf{Second prediction worth stating explicitly.} Low-temperature
LigandMPNN should \emph{saturate}: past some fold count it re-emits sequences
it has effectively already produced, so cumulative distinct solutions flatten
and further compute buys nothing. Report the saturation fold count as a number.
A visible plateau in the incumbent next to a still-climbing \cleo{} curve is
``filtering is not search'' rendered as a single line.}
\todo{ame --- run the $T$ sweep at matched fold budget; build the frontier panel}

\subsection*{Block 6 --- Rescue of backbones the benchmark cannot solve}

\plan{The strongest available result if it lands, and it is a clean one: the
near-miss class. Three backbones sit at \SIrange{1.61}{1.69}{\angstrom}
best-of-40 in the deposit --- \emph{zero} passing predictions, but only
$\sim$\SI{0.15}{\angstrom} from threshold. A pass on any of them is a design
the published pipeline cannot produce, with no reproduction confound, because
there is no positive result of theirs for us to fail to reproduce.

Present status: on the two pilot backbones our pipeline cannot reproduce,
\cleo{} recovers 4 passing designs each from 3200 folds where every baseline
arm returns zero. That is directionally right but too thin to build on, and it
is confounded by the reference-structure issue in Block 1. The near-miss
backbones are the experiment that makes this claim rather than hints at it ---
they have not yet been trained.}
\todo{ame --- train the three near-miss backbones; report rescue rate}
